\documentclass[10pt]{article}
\usepackage[margin=0.66in]{geometry}
\usepackage[T1]{fontenc}
\usepackage{helvet}\renewcommand{\familydefault}{\sfdefault}
\usepackage{booktabs,array,xcolor}
\definecolor{ink}{HTML}{1a1a1a}\definecolor{mut}{HTML}{736f6c}\definecolor{grn}{HTML}{034f46}
\color{ink}\setlength{\parindent}{0pt}\setlength{\parskip}{4pt}
\pagestyle{empty}
\begin{document}

{\large\textbf{One idea: ask the LLM how long a human would need, never for the probability}}\\
{\small\color{mut}GEPA forecaster project · full audit + two pre-registered live runs · 2026-09-02/03 · \$20.30 of the \$100 budget spent}

{\small\color{mut}The goal, restated. The risk models need P(model M succeeds at an attack step), given M's benchmark results. Real attack steps have no outcomes to score, so our stand-in for that jump is the held-out-benchmark test: train on 5 benchmarks, forecast CVEBench and CyberGym. A cell is one (model, task) pair. Brier is the error score, lower is better.}

\textbf{The trap we measured.} Every prompt that asked the LLM for the probability failed the same way on new benchmarks. Its ordering of tasks was fine, its numbers ran hot, and optimization made them hotter. Across all 8 GEPA prompts, a higher average forecast meant a worse test score, a perfect 1-for-1. GEPA inverted because its acceptance rule never scored a candidate on a held-out benchmark, so its edits pushed the one component that does not carry over.

\textbf{The idea.} The probability stops coming out of the LLM's mouth. Two parts:
\begin{enumerate}\setlength{\itemsep}{1pt}\setlength{\topsep}{2pt}
\item \textbf{A three-line prompt:} here is the task; how long would a skilled human need; answer low, mid, high in minutes. No model names, no evidence tables, no probabilities. This question the LLM answers well. Ranking a model's tasks by estimated time beats our curated difficulty labels at predicting which tasks it solves: 78\% vs 65\% pairwise-correct, $P=0.003$, on both test sets. And it works for any task, including an attack step with no ground truth.
\item \textbf{A 12-number calculator,} fit once on the training benchmarks: one ability number per model, one shared slope for how success falls with time. It converts (model, estimated minutes) into the probability. Audit the training data first. Leave-one-benchmark-out refits exposed InterCode-CTF as contaminated, solved 0.88--1.00 regardless of difficulty; with it removed, the calculator's average forecast matched the held-out base rate, 0.241 vs 0.240, with zero held-out labels.
\end{enumerate}

\begin{center}\small
\begin{tabular}{@{}lcc@{}}
\toprule
Brier & reserved test (1{,}033 cells) & fresh 27 tasks (297 cells, never prompted)\\
\midrule
seed prompt & 0.163 & 0.190\\
the 7 GEPA-optimized prompts & best ties seed, worst 0.230 & ---\\
calculator + curated difficulty labels & \textcolor{grn}{0.147} & --- \ (these tasks have no labels)\\
calculator + the dataset's LLM time estimates & 0.131--0.145 (unsettled) & \textcolor{grn}{0.155 \ ($t=-2.6$, pre-registered)}\\
calculator + our 3-line elicitation, end to end & \textcolor{grn}{0.151 \ (beats every GEPA prompt)} & 0.176 \ (right direction, underpowered)\\
\bottomrule
\end{tabular}
\end{center}

{\small We ran the whole pipeline live on 2026-09-03, pre-registered before launch: same LLM, temperature 0, one call per task, about \$12. On the reserved test it scored 0.151 vs the seed prompt's 0.163 and beat every GEPA prompt (task-level $t$ from $-3.7$ to $-8.4$). On the 27 fresh tasks: 0.176 vs 0.190, right direction, below the detectable-size bar we stated in advance. The dataset's own richer time estimates score 0.131--0.145 on the same test, so the three-line prompt has real headroom.}

\textbf{The honest trade.} In distribution, where good curated labels exist, the elicited times lose: sealed 0.142 vs the seed prompt's 0.131. Rule: use labels where you have them, elicit where you do not; new benchmarks and attack steps are exactly where you do not. The forecast scale still runs high on new benchmarks, average 0.343 vs realized 0.240; the ranking does the work. One correction number is held in reserve for the first real outcomes. On the studied test it repaired the scale after 50 outcomes, 0.145 to 0.135; on the fresh run the scale needed no fix. This page replaces the earlier one: that page's curve-plus-seed blend added nothing on fresh cells and is retired.

\textbf{Limits.} Two held-out benchmarks, mostly CyberGym. Only the fresh run and the live pipeline were pre-registered; the rest was computed after the test labels had been studied. The time-ranking result was found independently in the dataset dossier first. The losing curated labels are mostly expert guesses themselves. A single difficulty number per task ignores model quirks such as refusals; no measured prompt captured those either. All of it is a proxy for real attack steps, which stay unmeasurable.

\textbf{Next, if the team wants it.} Point GEPA at the three-line time prompt, never again at a probability prompt, and make it hold out one training benchmark per fold when accepting candidates, about \$150. In deployment, the first 20 to 50 real outcomes decide whether the correction number is needed.

\vspace{2pt}
{\footnotesize\color{mut}Every number reproduced by at least two independent derivations from raw logs. Pre-registrations committed before launch: gepa \texttt{runs/holdout27\_prereg.md}, \texttt{runs/darm\_prereg.md}. Raw cells on the results branch. Errata appended to \texttt{sweep\_and\_test\_2026-09-01.md}. One dataset, one harness, one 11-model panel.}

\end{document}
